\documentclass[11pt,a4paper]{article}

% ===================== PREÁMBULO =====================
\usepackage[utf8]{inputenc}
\usepackage[T1]{fontenc}
\usepackage{lmodern}
% Nombres en español sin depender de babel-spanish (no disponible en este entorno)
\renewcommand{\contentsname}{Índice}
\renewcommand{\refname}{Referencias}
\renewcommand{\tablename}{Tabla}
\renewcommand{\figurename}{Figura}
\renewcommand{\abstractname}{Resumen}
\usepackage{amsmath,amssymb,mathtools}
% --- Mini-sistema de unidades (sustituye a siunitx, no disponible en el entorno) ---
\providecommand{\sisetup}[1]{}
\newcommand{\si}[1]{\ensuremath{#1}}
\newcommand{\SI}[2]{\ensuremath{#1\,#2}}
\newcommand{\watt}{\mathrm{W}}
\newcommand{\meter}{\mathrm{m}}
\newcommand{\second}{\mathrm{s}}
\newcommand{\gram}{\mathrm{g}}
\newcommand{\joule}{\mathrm{J}}
\newcommand{\liter}{\mathrm{L}}
\newcommand{\lux}{\mathrm{lx}}
\newcommand{\lumen}{\mathrm{lm}}
\newcommand{\mol}{\mathrm{mol}}
\newcommand{\steradian}{\mathrm{sr}}
\renewcommand{\deg}{{}^{\circ}}
\newcommand{\degree}{{}^{\circ}}
\newcommand{\per}{/}
\renewcommand{\square}[1]{#1^{2}}
\newcommand{\cubic}[1]{#1^{3}}
\newcommand{\milli}[1]{\mathrm{m}#1}
\newcommand{\centi}[1]{\mathrm{c}#1}
\newcommand{\micro}[1]{\upmu #1}
\newcommand{\nano}[1]{\mathrm{n}#1}
\newcommand{\arcmin}{\mathrm{arcmin}}
\providecommand{\upmu}{\mu}
\usepackage{booktabs}
\usepackage{graphicx}
\usepackage{geometry}
\usepackage{xcolor}
\usepackage{enumitem}
\usepackage{longtable}
\usepackage[hidelinks]{hyperref}
\usepackage{caption}

\geometry{margin=2.5cm}
\hypersetup{pdftitle={Documentacion fisica del simulador light-cage-sim},
            pdfauthor={EVOLUX}}

\newcommand{\code}[1]{\texttt{#1}}
\newcommand{\fn}[1]{\textcolor{teal}{\texttt{#1}}}
% Caja de trazabilidad al código fuente
\newcommand{\src}[1]{\par\vspace{2pt}\noindent\fcolorbox{gray!40}{gray!8}{%
  \parbox{0.985\linewidth}{\footnotesize\textbf{Código:} #1}}\par\vspace{4pt}}
% Nota de modelado / advertencia
\newcommand{\nota}[1]{\par\vspace{2pt}\noindent\fcolorbox{orange!50}{orange!8}{%
  \parbox{0.985\linewidth}{\footnotesize\textbf{Nota de modelado:} #1}}\par\vspace{4pt}}

\title{\textbf{Documentación física del simulador \code{light-cage-sim}}\\[4pt]
\large Trazado de rayos Monte Carlo para irradiancia submarina en jaulas y estanques de acuicultura}
\author{EVOLUX}
\date{\today}

\begin{document}
\maketitle

\begin{abstract}
\noindent Este documento describe, con nivel técnico de ingeniería, los modelos
físicos implementados en el simulador \code{light-cage-sim}: un trazador de rayos
estocástico que estima la distribución espacial y espectral de la luz emitida por
luminarias sumergidas o aéreas en recintos de acuicultura (jaulas marinas y
estanques RAS). Se cubren el marco radiométrico y fotométrico, la lectura de la
fotometría de la fuente (IES~TM-33-18 e IES~LM-63), el muestreo espectral, la
refracción y reflexión de Fresnel en la interfaz aire--agua, el modelo bio-óptico
de cuatro componentes y su relación con las propiedades aparentes del agua, el
transporte radiativo por la ley de Beer--Lambert y por Monte Carlo volumétrico con
función de fase de Henyey--Greenstein, y la conversión final a unidades
fotométricas y fotosintéticas. Cada ecuación se acompaña de su trazabilidad al
código fuente y se distingue explícitamente entre física fundamental y decisiones
de modelado del simulador. La bibliografía referencia las fuentes primarias de
cada submodelo.
\end{abstract}

\tableofcontents
\bigskip

% =====================================================
\section{Alcance y arquitectura del modelo}
% =====================================================

El simulador resuelve, de forma numérica y estocástica, la propagación de la luz
desde una o más luminarias hacia planos horizontales de observación situados a
profundidades objetivo dentro de un recinto acuático. El recinto admite dos
geometrías y dos convenciones verticales: el \emph{estanque}, con cota medida desde
el nivel de piso y superficie del agua en $z=z_\text{int}$; y la \emph{jaula}, con
profundidad medida desde la superficie ($z=0$) y suelo en $z=-z_\text{env}$. La
sección transversal puede ser circular o rectangular.

\src{\fn{SimulationEngine.run} (\code{simulation\_engine.py}, líneas 196--640).
Convención vertical: líneas 219--224 y 268; geometría: líneas 200--211.}

El flujo de cálculo consta de seis etapas encadenadas, todas vectorizadas con
\code{NumPy}:
\begin{enumerate}[itemsep=2pt]
  \item Lectura de la distribución de intensidad de la luminaria
        (\S\ref{sec:fotometria}).
  \item Generación cuasi-uniforme de direcciones de rayo sobre la esfera y
        normalización a la potencia radiante objetivo (\S\ref{sec:muestreo}).
  \item Asignación espectral por muestreo de la distribución de potencia
        (\S\ref{sec:espectral}).
  \item Refracción y transmitancia de Fresnel en la interfaz aire--agua, cuando la
        fuente es aérea (\S\ref{sec:fresnel}).
  \item Transporte en el medio acuático, por uno de dos modos: Beer--Lambert
        determinista (\S\ref{sec:beer}) o Monte Carlo con dispersión
        (\S\ref{sec:montecarlo}). Las propiedades ópticas del agua provienen del
        modelo bio-óptico (\S\ref{sec:biooptico}).
  \item Agregación de los cruces de rayo en una malla y conversión a irradiancia y
        unidades derivadas (\S\ref{sec:agregacion}).
\end{enumerate}

% =====================================================
\section{Marco radiométrico y fotométrico}\label{sec:radiometria}
% =====================================================

La magnitud primaria que transporta cada rayo es flujo radiante $\Phi$
(\si{\watt}). La intensidad radiante $I(\theta,\phi)$ (\si{\watt\per\steradian})
describe la potencia por unidad de ángulo sólido emitida por la fuente en la
dirección $(\theta,\phi)$. La potencia radiante total de la fuente es

\begin{equation}
\Phi_\text{rad} \;=\; \oint_{4\pi} I(\theta,\phi)\,\mathrm{d}\Omega
\;=\; \int_0^{2\pi}\!\!\int_0^{\pi} I(\theta,\phi)\,\sin\theta\,\mathrm{d}\theta\,\mathrm{d}\phi .
\label{eq:flujo}
\end{equation}

La cantidad de interés en cada plano de observación es la \emph{irradiancia} $E$
(\si{\watt\per\square\meter}), potencia por unidad de área que atraviesa el plano.
El simulador soporta tres definiciones operativas del sensor:

\begin{description}[style=nextline,leftmargin=1.4em]
  \item[Escalar / planar.] Se contabiliza el flujo que cruza el plano horizontal,
  con signo resuelto por el sentido del rayo. La cantidad reportada equivale a
  $|E_d| + |E_u|$ (descendente más ascendente), que en ausencia de dispersión se
  reduce a la irradiancia direccional clásica.
  \item[Descendente (\code{downwelling}).] Representa un colector plano horizontal
  con su cara activa mirando hacia arriba. Sólo contabiliza rayos descendentes
  ($D_z<0$), por lo que entrega $E_d$ y excluye el flujo ascendente que incidiría
  sobre la cara posterior del instrumento.
  \item[Pineal.] Modela un fotorreceptor que mira hacia arriba con una respuesta
  angular coseno-ponderada. Cada rayo descendente que cruza el plano se pondera por
  \begin{equation}
  w_\text{pineal}(\mu) = \tfrac{1}{2}\,\bigl(1+\cos\mu\bigr)\quad\text{si }\cos\mu \ge \cos\mu_\text{max},\qquad 0 \text{ en otro caso},
  \label{eq:pineal}
  \end{equation}
  donde $\mu$ es el ángulo respecto al cenit del sensor y $\mu_\text{max}$ un ángulo
  de corte (por defecto \SI{85}{\degree}). El factor $\tfrac12$ se aplica sólo si la
  normalización está activa.
\end{description}

\src{Tipos de irradiancia: parámetro \code{irradiance\_type}
(\code{simulation\_engine.py}, líneas 247--251). Ponderación pineal: líneas
437--443 (modo Beer--Lambert) y 555--559 (modo Monte Carlo).}

\nota{La etiqueta ``pineal'' designa una respuesta angular tipo coseno con corte,
no una curva de acción espectral del órgano pineal. Es una elección de modelado del
simulador; no debe interpretarse como una sensibilidad fotobiológica calibrada.}

\subsection{Conversión a unidades fotométricas y fotosintéticas}\label{sec:unidades}

Tras la simulación radiométrica, la energía espectral de cada impacto se convierte
a dos escalas de uso práctico. La iluminancia en lux se obtiene ponderando por la
eficiencia luminosa espectral fotópica $V(\lambda)$ y la constante de eficacia
luminosa máxima $K_m = \SI{683}{\lumen\per\watt}$:

\begin{equation}
E_v \;=\; K_m \int V(\lambda)\, E_\lambda(\lambda)\,\mathrm{d}\lambda .
\label{eq:lux}
\end{equation}

El simulador evalúa $V(\lambda)$ mediante el ajuste analítico gaussiano doble de
Wyman--Sloan--Shirley con $\lambda$ en \si{\micro\meter}:

\begin{equation}
V(\lambda) = 1{,}019\,e^{-285{,}4\,(\lambda-0{,}559)^2}
            \;-\; 0{,}092\,e^{-1250\,(\lambda-0{,}450)^2},
\label{eq:vlambda}
\end{equation}

acotado a $[0,1]$. El flujo de fotones fotosintéticos (PPFD, en
\si{\micro\mol\per\square\meter\per\second}) se calcula sólo en la banda PAR
($400\text{--}\SI{700}{\nano\meter}$), donde el número de fotones por julio es
proporcional a $\lambda$:

\begin{equation}
\text{PPFD} \;=\; \int_{400}^{700} \frac{\lambda\,[\si{\nano\meter}]}{119{,}626}\, E_\lambda(\lambda)\,\mathrm{d}\lambda .
\label{eq:ppfd}
\end{equation}

La constante $119{,}626$ equivale a $h c N_A \times 10^{6}$ en unidades de
$\si{\nano\meter}\cdot\si{\joule}\cdot\si{\micro\mol}^{-1}$ y proviene de la energía de un fotón
$\varepsilon = hc/\lambda$ junto con el número de Avogadro.

\src{Conversión por impacto: \code{app\_sim.py}, líneas 453--461. Factor de
escala medio por capa \code{f\_lux}, \code{f\_ppfd} ponderado por la energía total
de los impactos.}

% =====================================================
\section{Fotometría de la fuente}\label{sec:fotometria}
% =====================================================

La distribución angular de emisión de cada luminaria se lee de archivos
fotométricos estándar. El simulador admite dos familias:

\begin{itemize}[itemsep=2pt]
  \item \textbf{IES TM-33-18} (XML). Se extraen dos campos angulares: intensidad
  luminosa (\code{LuminousIntensity}, en candelas) e intensidad radiante
  (\code{RadiantIntensity}, en \si{\watt\per\steradian}), tabuladas por ángulo
  horizontal $h$ y vertical $v$. Adicionalmente se leen la potencia eléctrica de
  entrada (\code{InputWattage}), el flujo radiante total (\code{RadiantFlux}) y la
  distribución espectral de potencia (\code{EmitterSpectral}).
  \item \textbf{IES LM-63} (\code{.ies}). Formato fotométrico clásico en candelas,
  con simetría angular reconstruida según el rango de ángulos horizontales
  (cuadrante, semiplano o completo).
\end{itemize}

\src{\fn{TM33Parser} e \fn{IESParser} en \code{parsers.py} (líneas 6--202).
Reconstrucción de simetría azimutal IES: líneas 172--187.}

La intensidad en una dirección arbitraria $\hat{\mathbf{d}}$ se obtiene por
interpolación bilineal sobre la rejilla angular. La conversión de la dirección
cartesiana del rayo a coordenadas fotométricas usa

\begin{equation}
v = \arccos(-d_z), \qquad h = \operatorname{atan2}(d_y, d_x)\ \bmod 2\pi,
\label{eq:angulos}
\end{equation}

de modo que $v$ es el ángulo polar medido desde el nadir de la luminaria.

\src{\fn{TM33Parser.get\_intensity}, \code{parsers.py}, líneas 102--111;
interpolador \code{RegularGridInterpolator} con extrapolación nula
(\code{fill\_value=0}).}

\subsection{Normalización a potencia radiante}

La rejilla de intensidad describe la \emph{forma} angular de emisión, no su escala
absoluta. El simulador renormaliza la intensidad radiante muestreada para que su
integral angular reproduzca la potencia óptica objetivo
$\Phi_\text{obj}=P_\text{elec}\,\eta$, donde $P_\text{elec}$ es la potencia
eléctrica y $\eta$ la eficiencia óptica del equipo. Con $N$ rayos generados sobre la
esfera, la cuadratura Monte Carlo de \eqref{eq:flujo} es

\begin{equation}
\Phi_\text{actual} \;\approx\; \frac{4\pi}{N}\sum_{k=1}^{N} I(\hat{\mathbf{d}}_k),
\qquad
I \leftarrow I\,\frac{\Phi_\text{obj}}{\Phi_\text{actual}} .
\label{eq:norm}
\end{equation}

El flujo que transporta cada rayo incorpora además el factor de \emph{dimming}
$\delta\in[0,1]$:
\begin{equation}
\Phi_k \;=\; I(\hat{\mathbf{d}}_k)\,\frac{4\pi}{N}\,\delta .
\label{eq:fluxray}
\end{equation}

\src{Normalización y flujo por rayo: \code{simulation\_engine.py}, líneas
284--298.}

% =====================================================
\section{Muestreo direccional}\label{sec:muestreo}
% =====================================================

Las direcciones de emisión se distribuyen de forma cuasi-uniforme sobre la esfera
mediante la espiral de Fibonacci, que evita el agrupamiento polar del muestreo en
latitud--longitud y reduce la varianza respecto al muestreo puramente aleatorio.
Para $k=0,\dots,N-1$ y $i_k=k+\tfrac12$:

\begin{align}
\phi_k &= \arccos\!\Bigl(1 - \tfrac{2 i_k}{N}\Bigr), &
\theta_k &= \pi\,(1+\sqrt5)\, i_k, \\[2pt]
\hat{\mathbf{d}}_k &= \bigl(\sin\phi_k\cos\theta_k,\ \sin\phi_k\sin\theta_k,\ \cos\phi_k\bigr). &&
\end{align}

\src{\code{simulation\_engine.py}, líneas 276--280. El término $(1+\sqrt5)$
corresponde al ángulo áureo.}

\subsection{Rotación de la luminaria}

La orientación del equipo se aplica con una rotación rígida $R = R_z R_y R_x$
construida con las matrices de rotación elementales en torno a los ejes
$x$, $y$, $z$ por los ángulos $(\rho_x,\rho_y,\rho_z)$. Los rayos en marco local se
transforman al marco global como $\hat{\mathbf{d}}_\text{glob} = R\,\hat{\mathbf{d}}_\text{loc}$.

\src{\fn{rotate\_3d}, \code{simulation\_engine.py}, líneas 28--44; aplicación en
líneas 300--302.}

% =====================================================
\section{Muestreo espectral}\label{sec:espectral}
% =====================================================

A cada rayo se le asigna una longitud de onda $\lambda$ muestreada de la
distribución espectral de potencia de la fuente $\{(\lambda_j, p_j)\}$. La densidad
de probabilidad se supone lineal a trozos, por lo que la masa de cada segmento
entre $\lambda_j$ y $\lambda_{j+1}$ es el área del trapecio

\begin{equation}
A_j \;=\; \tfrac12\,(p_j + p_{j+1})\,(\lambda_{j+1}-\lambda_j),
\qquad
\text{CDF}_m \;=\; \frac{\sum_{j<m} A_j}{\sum_j A_j}.
\label{eq:cdf}
\end{equation}

El muestreo se realiza invirtiendo la función de distribución acumulada (CDF) sobre
un número uniforme $\xi\sim U(0,1)$ por interpolación lineal. Esta construcción es
más fiel que una suma acumulada simple cuando la rejilla espectral no es uniforme.
Si la luminaria carece de datos espectrales (caso típico de archivos IES), se asume
un espectro plano en $\{400,500,600,700\}\,\si{\nano\meter}$.

\src{\fn{sample\_wavelength}, \code{simulation\_engine.py}, líneas 97--112;
fallback espectral plano en líneas 306--312.}

% =====================================================
\section{Interfaz aire--agua: refracción y Fresnel}\label{sec:fresnel}
% =====================================================

Cuando la luminaria está sobre la superficie del agua (sólo en geometría de
estanque, con $z_\text{lampara} > z_\text{int}$), cada rayo descendente que alcanza
la interfaz se refracta según la ley de Snell. Sean $n_1$ el índice del aire
(por defecto $1{,}0$) y $n_2$ el del agua (por defecto $1{,}333$), y $\theta_i$ el
ángulo de incidencia respecto a la normal vertical. La refracción cumple

\begin{equation}
n_1 \sin\theta_i = n_2 \sin\theta_t,
\qquad
\sin^2\theta_t = \Bigl(\tfrac{n_1}{n_2}\Bigr)^2 (1-\cos^2\theta_i).
\label{eq:snell}
\end{equation}

El vector transmitido se obtiene en forma vectorial:
\begin{equation}
\hat{\mathbf{t}} \;=\; \frac{n_1}{n_2}\,\hat{\mathbf{d}}
\;+\;\Bigl(\tfrac{n_1}{n_2}\cos\theta_i - \cos\theta_t\Bigr)\,\hat{\mathbf{n}},
\label{eq:tvec}
\end{equation}
con $\hat{\mathbf{n}}=(0,0,1)$. La fracción de potencia transmitida usa el
coeficiente de Fresnel para luz no polarizada (promedio de las componentes
$s$ y $p$):

\begin{align}
R_s &= \left(\frac{n_1\cos\theta_i - n_2\cos\theta_t}{n_1\cos\theta_i + n_2\cos\theta_t}\right)^{\!2}, &
R_p &= \left(\frac{n_1\cos\theta_t - n_2\cos\theta_i}{n_1\cos\theta_t + n_2\cos\theta_i}\right)^{\!2}, \\[4pt]
T &= 1 - \tfrac12 (R_s + R_p). &&
\label{eq:fresnel}
\end{align}

El flujo del rayo transmitido se multiplica por $T$. En el sentido aire$\to$agua
($n_1<n_2$) no existe reflexión interna total, pero el código conserva la prueba
$\sin^2\theta_t\le1$ por robustez numérica.

\src{\fn{fresnel\_transmission}, \code{simulation\_engine.py}, líneas 15--20;
refracción aire$\to$agua en líneas 363--392. El comentario ``CORRECCIÓN \#1''
documenta la eliminación de un factor empírico $0{,}98$ redundante, ya que Fresnel
conserva la potencia por sí mismo.}

En el modo de dispersión, la interfaz también actúa para rayos ascendentes dentro
del agua (sentido agua$\to$aire, $n_2>n_1$). Allí sí aparece reflexión interna
total cuando $\sin^2\theta_t\ge1$, en cuyo caso la reflectancia es $R=1$; en otro
caso $R = 1-T$ y el rayo se refleja especularmente respecto a la normal vertical.

\src{Reflexión en superficie con TIR (modo Monte Carlo):
\code{simulation\_engine.py}, líneas 594--615.}

% =====================================================
\section{Modelo bio-óptico de cuatro componentes}\label{sec:biooptico}
% =====================================================

Las propiedades ópticas inherentes (IOP) del agua se construyen como suma de cuatro
contribuyentes: agua pura, materia orgánica disuelta coloreada (CDOM), fitoplancton
(vía clorofila-a) y sólidos suspendidos totales (TSS). El coeficiente de absorción
total $a(\lambda)$ y el coeficiente de dispersión total $b(\lambda)$ son

\begin{align}
a(\lambda) &= a_w(\lambda) \;+\; a_\text{CDOM}(440)\,e^{-S_\text{CDOM}(\lambda-440)}
              \;+\; a_\phi^{*}(\lambda)\,\text{Chl}, \label{eq:abs}\\[3pt]
b(\lambda) &= b^{*}(\lambda)\,\text{TSS}. \label{eq:scat}
\end{align}

Aquí $a_w$ es la absorción del agua pura; el término CDOM decae exponencialmente
desde su valor de referencia a \SI{440}{\nano\meter} con pendiente espectral
$S_\text{CDOM}$ (por defecto \SI{0.015}{\per\nano\meter}); $a_\phi^{*}$ es la
absorción específica de la clorofila; y $b^{*}$ la dispersión específica por unidad
de masa de TSS. El coeficiente de atenuación de haz es $c(\lambda)=a(\lambda)+b(\lambda)$
y el albedo de dispersión simple $\omega(\lambda)=b(\lambda)/c(\lambda)$.

\src{\fn{bio\_optical\_iop}, \code{simulation\_engine.py}, líneas 132--150.
Activación con \code{optics.mc\_input\_type='bio'}: líneas 323--330.}

Los espectros de referencia están tabulados en una rejilla de
$400$--$\SI{700}{\nano\meter}$ e interpolados linealmente. La
Tabla~\ref{tab:iop} reproduce los valores codificados.

\begin{table}[h]
\centering
\caption{Coeficientes específicos de referencia del modelo bio-óptico
(\code{simulation\_engine.py}, líneas 119--129).}
\label{tab:iop}
\small
\begin{tabular}{lccccccc}
\toprule
$\lambda$ [\si{\nano\meter}] & 400 & 450 & 500 & 550 & 600 & 650 & 700 \\
\midrule
$a_w$ [\si{\per\meter}] & 0,018 & 0,015 & 0,026 & 0,064 & 0,245 & 0,349 & 0,624 \\
$b^{*}$ [\si{\square\meter\per\gram}] & 0,50 & 0,42 & 0,35 & 0,31 & 0,28 & 0,25 & 0,22 \\
$a_\phi^{*}$ [\si{\square\meter\per\milli\gram}] & 0,022 & 0,038 & 0,012 & 0,005 & 0,005 & 0,018 & 0,008 \\
\bottomrule
\end{tabular}
\end{table}

\nota{$a_w$ deriva de Pope \& Fry (1997)~\cite{popefry} y Smith \& Baker
(1981)~\cite{smithbaker}, redondeados; $a_\phi^{*}$ es un promedio de Bricaud et
al.\ (1995, 1998)~\cite{bricaud}; $b^{*}$ es un perfil representativo de partículas
minerales coherente con Babin et al.\ (2003)~\cite{babin}. Las concentraciones de
entrada por defecto del modo \code{bio} (TSS, CDOM, Chl) son valores conservadores y
no constituyen mediciones de campo.}

\subsection{Fundamento físico de cada componente}\label{sec:componentes}

La descomposición aditiva de las ecuaciones \eqref{eq:abs}--\eqref{eq:scat} es el
modelo bio-óptico canónico de aguas Caso~1/Caso~2, en el que la absorción total se
escribe como la suma de los aportes independientes del agua, el fitoplancton, el
CDOM y los detritos/minerales~\cite{ioccg,morel,babin}. El simulador adopta sus
formas funcionales estándar:

\begin{itemize}[itemsep=3pt]
  \item \textbf{Agua pura.} $a_w(\lambda)$ es una propiedad intrínseca del medio,
  prácticamente invariante salvo dependencias menores con temperatura y salinidad
  en el infrarrojo cercano~\cite{popefry,smithbaker}. Domina la absorción por encima
  de \SI{600}{\nano\meter}, lo que explica la atenuación selectiva del rojo en
  profundidad.
  \item \textbf{CDOM.} La absorción de la materia orgánica disuelta coloreada decae
  cuasi-exponencialmente con la longitud de onda~\cite{bricaud_cdom,kirk}, lo que
  justifica la parametrización de un único punto de anclaje a \SI{440}{\nano\meter}
  y la pendiente espectral $S_\text{CDOM}$. El rango $0{,}014$--$\SI{0.018}{\per\nano\meter}$
  empleado por defecto es el típico de aguas costeras y de fiordo~\cite{bricaud_cdom}.
  \item \textbf{Fitoplancton.} La absorción específica $a_\phi^{*}(\lambda)$ presenta
  los máximos característicos de la clorofila-a a $\sim$\SI{440}{} y
  $\sim$\SI{675}{\nano\meter}; su forma y amplitud varían con el ``efecto
  paquete''~\cite{bricaud}, que el modelo no resuelve (usa un $a_\phi^{*}$ fijo).
  \item \textbf{Sólidos suspendidos (dispersión).} La dispersión total se modela
  proporcional a la concentración de TSS mediante una sección eficaz específica
  $b^{*}(\lambda)$ que decrece suavemente con $\lambda$, consistente con partículas
  minerales de distribución de tamaños tipo Junge~\cite{babin}.
\end{itemize}

El modelo asume implícitamente que la retrodispersión es una fracción fija de la
dispersión total a través del factor de asimetría $g$ (\S\ref{sec:montecarlo}); no
parametriza un coeficiente de retrodispersión $b_b(\lambda)$ independiente. Esta es
una simplificación deliberada coherente con el uso de una única función de fase de
Henyey--Greenstein.

\subsection{De propiedades inherentes a aparentes}

Las propiedades ópticas inherentes (IOP) —$a$, $b$, $c$— son independientes del
campo de luz, mientras que las propiedades ópticas aparentes (AOP) —como $K_d$—
dependen también de la geometría de iluminación~\cite{mobley,kirk}. El puente entre
ambas es necesario porque la teledetección y la práctica de campo entregan
habitualmente $K_d$, mientras que el trazador de rayos necesita $c$ y $\omega$.

El punto de partida formal es la ley de Gershun, que relaciona la divergencia del
vector irradiancia con la absorción neta. Para un campo cuasi-colimado descendente
con coseno medio $\bar\mu_d$, la teoría de transferencia radiativa de
Kirk~\cite{kirk,kirk1984} conduce a una forma de dos términos en la que la
atenuación difusa se debe a la absorción más una contribución de la dispersión
ponderada por su asimetría:

\begin{equation}
K_d \;\approx\; \frac{a + (1-g)\,b}{\bar\mu_d},
\label{eq:kdfromiop}
\end{equation}

donde $g$ es el factor de asimetría de la dispersión y $\bar\mu_d$ el coseno medio
del campo descendente (ambos por defecto $0{,}85$). El término $(1-g)\,b$ aproxima
la dispersión efectiva hacia adelante perdida del haz directo y desempeña el papel
de la retrodispersión en formulaciones semianalíticas como la de Lee et
al.~\cite{lee2005}. La relación inversa, que entrega un trío $(c,a,b)$ consistente a
partir de un $K_d$ dado suponiendo un albedo $\omega=b/c$ y un $g$ típicos, se
obtiene sustituyendo $a=c(1-\omega)$ y $b=\omega c$ en \eqref{eq:kdfromiop}:

\begin{equation}
K_d \approx \frac{c(1-\omega) + (1-g)\,\omega c}{\bar\mu_d}
\;\Longrightarrow\;
c = \frac{K_d\,\bar\mu_d}{1-\omega g},
\qquad b = \omega\, c,
\qquad a = c-b.
\label{eq:cfromkd}
\end{equation}

\src{\fn{kd\_from\_iop} (líneas 153--159) y \fn{c\_from\_kd} (líneas 162--173) en
\code{simulation\_engine.py}. La derivación algebraica de $c$ está documentada en el
\emph{docstring} de \fn{c\_from\_kd} (líneas 165--167).}

\nota{Las ecuaciones \eqref{eq:kdfromiop}--\eqref{eq:cfromkd} son aproximaciones de
cierre, no identidades exactas. $\bar\mu_d$ no es constante: aumenta con la
profundidad a medida que el campo se difumina, por lo que tomarlo fijo en $0{,}85$
introduce un sesgo creciente en aguas muy dispersoras. La validez es mejor en aguas
de oligotróficas a mesotróficas y se degrada con cargas de partículas muy altas o
geometrías de iluminación fuertemente direccionales~\cite{lee2005,ioccg}.}

\subsubsection*{Cierre semianalítico de Lee et al.\ (2005)}\label{sec:leekd}

Como alternativa de mayor fidelidad al cierre tipo Kirk, el simulador ofrece el
modelo semianalítico de Lee, Du \& Arnone (2005)~\cite{lee2005}
(\code{kd\_closure='lee2005'}), que usa la retrodispersión $b_b$ de forma
\emph{explícita} y la geometría de iluminación a través del ángulo cenital solar en
aire $\theta_a$:

\begin{equation}
K_d = (1 + 0{,}005\,\theta_a)\,a \;+\; 4{,}18\,\bigl(1 - 0{,}52\,e^{-10{,}8\,a}\bigr)\,b_b,
\label{eq:kdlee}
\end{equation}

con $\theta_a$ en grados (por defecto $30^\circ$, nominal, al tratarse de una fuente
artificial). La retrodispersión se toma de la fase activa: $b_b = B\,b$, con
$B=B(g)$ para Henyey--Greenstein \eqref{eq:hgbackscatter} o $B=b_b/b$ objetivo para
Fournier--Forand \eqref{eq:ffbb}. Esto enlaza de forma coherente la elección de fase
(\S\ref{sec:ff}) con el cálculo de $K_d$ y, por ende, con la profundidad de Secchi.

\src{\fn{kd\_lee2005} en \code{simulation\_engine.py}; selección de cierre y cálculo
de $B$ activo en \code{app\_sim.py} (función \fn{\_kd\_from\_iop\_active}).}

\nota{Frente al cierre de Kirk, \eqref{eq:kdlee} entrega típicamente un $K_d$ algo
menor en aguas dispersoras (p.\ ej.\ $\sim$5--10\,\% a $440$--$\SI{650}{\nano\meter}$
para una carga costera moderada), porque pondera $b_b$ —no la dispersión total— y
captura la saturación del término de dispersión con la absorción.}

\subsection{Modo RAS (Bårdsnes 2020): formas calibradas, magnitud
calibrable}\label{sec:ras}

El simulador expone un modo \code{scattering\,$\to$\,ras\_bardsnes} orientado a agua
de sistemas de recirculación (RAS). A diferencia del modo marino, la atenuación
\emph{crece hacia el azul} —inverso al océano— por la combinación de CDOM concentrado
y micropartículas orgánicas finas. La estructura es de cuatro componentes, con la
atenuación particulada representada como ley de potencia:

\begin{align}
c_p(\lambda) &= b^{*}_{550}\,[\mathrm{TSS}]\left(\frac{\lambda}{550}\right)^{-\eta_p},
  \label{eq:rascp}\\
b(\lambda) &= \omega_p\,c_p(\lambda), \qquad
a_p(\lambda) = (1-\omega_p)\,c_p(\lambda), \label{eq:rasbap}\\
a(\lambda) &= a_w(\lambda) + a_{440}\,e^{-S_{\mathrm{CDOM}}(\lambda-440)}
  + a^{*}_{phy}(\lambda)\,[\mathrm{Chl}] + a_p(\lambda). \label{eq:rasa}
\end{align}

\paragraph{Qué proviene del trabajo original y qué no.} De Bårdsnes (2020) se toman
las \emph{formas} espectrales: la pendiente particulada $\eta_p \approx 1{,}8$
—ajustada a la columna TSS de la Tabla 4.1, más pronunciada que el $\sim$1,5 marino
porque las partículas de RAS son más finas— y la pendiente de absorción del CDOM
$S_{\mathrm{CDOM}} \approx 0{,}0141\;\mathrm{nm^{-1}}$ (ajuste 400--500\,nm, columna
DOC). También la conversión de turbidez a sólidos suspendidos medida en tanque,
$\mathrm{TSS} = 3{,}0411\,\mathrm{NTU} - 0{,}376$ ($R^2 = 0{,}86$, Fig.\ 3.2A).

La \emph{magnitud absoluta} no es transferible: la medición del trabajo original
tiene re-entrada de luz por las paredes del tanque y el propio autor indica que la
absorbancia «aparece menor que en la realidad». Por eso $b^{*}_{550}$ (atenuación
específica particulada) y $\omega_p$ (albedo de dispersión simple) quedan expuestos
como parámetros \textbf{calibrables por instalación} en la interfaz, con valores por
defecto ($0{,}31\;\mathrm{m^2/g}$ y $0{,}90$) elegidos para preservar continuidad de
magnitud con el modo marino, no por ser universales.

\nota{Antes de usar esta ruta para dimensionar, calibre $b^{*}_{550}$ y $\omega_p$
con una medida óptica del propio sistema: $c(\lambda)$, $K_d(\lambda)$ o
transmitancia espectral. Sin esa calibración el modo es útil para explorar el
\emph{color} y la forma espectral del campo de luz en RAS, no su magnitud.}

\src{\fn{bio\_optical\_iop\_ras\_bardsnes} y \fn{ras\_tss\_from\_turbidity} en
\code{simulation\_engine.py}; selección del modo en \code{app\_sim.py}
(\code{mc\_input\_type == 'ras\_bardsnes'}).}

% =====================================================
\section{Transporte radiativo I: Beer--Lambert}\label{sec:beer}
% =====================================================

En los modos \code{kd\_fijo} y \code{kd\_espectral} la atenuación es determinista y
sin dispersión múltiple. El simulador distingue cuidadosamente dos interpretaciones
del coeficiente de extinción, controladas por \code{atten\_coef\_type}:

\begin{itemize}[itemsep=3pt]
  \item \textbf{Coeficiente de haz $c$.} La pérdida se aplica a lo largo de la
  longitud real del camino del rayo $d_\text{path}$:
  \begin{equation}
  \Phi(d_\text{path}) = \Phi_0\, e^{-c\, d_\text{path}}.
  \label{eq:beam}
  \end{equation}
  \item \textbf{Coeficiente difuso $K_d$.} Definido para la irradiancia descendente
  frente a la profundidad, se aplica al desplazamiento \emph{vertical}
  $\Delta z = |z_\text{hit}-z_0|$, reproduciendo en el agregado la ley clásica
  \begin{equation}
  E_d(z) = E_d(0)\, e^{-K_d\, z}.
  \label{eq:kd}
  \end{equation}
\end{itemize}

En el modo \code{kd\_espectral}, $K_d(\lambda)$ (o $c(\lambda)$) se interpola por
longitud de onda a partir de una tabla provista. Cada rayo cruza el plano de
observación a lo sumo una vez; el signo del parámetro de intersección $t$ resuelve
de forma natural tanto los planos situados por debajo como por encima de la fuente,
de modo que la cantidad acumulada equivale a $|E_d|+|E_u|$.

\src{Modo Beer--Lambert: \code{simulation\_engine.py}, líneas 400--449. Distinción
$K_d$ vs $c$ (``CORRECCIÓN \#7''): líneas 427--435.}

% =====================================================
\section{Transporte radiativo II: Monte Carlo con dispersión}\label{sec:montecarlo}
% =====================================================

El modo \code{scattering} resuelve la ecuación de transferencia radiativa por el
método de Monte Carlo, propagando paquetes de energía (``fotones ponderados'') a
través de eventos sucesivos de dispersión y de interacción con las fronteras del
recinto. Cada paquete porta un peso $W$ (su flujo) que se atenúa multiplicativamente
en cada interacción.

\subsection{Distancia libre y clasificación de eventos}

La distancia hasta el próximo evento de dispersión volumétrica se muestrea de la
distribución exponencial de camino libre, característica de un medio con coeficiente
de atenuación $c$:

\begin{equation}
t_\text{scat} = -\frac{\ln \xi}{c}, \qquad \xi \sim U(0,1).
\label{eq:freepath}
\end{equation}

En paralelo se calculan las distancias geométricas a las tres fronteras: la pared
lateral $t_\text{wall}$ (intersección rayo--cilindro o rayo--caja), el suelo
$t_\text{floor}$ y la superficie $t_\text{surf}$. El evento que ocurre es el de
menor distancia,
\begin{equation}
t_\text{event} = \min\{t_\text{wall}, t_\text{floor}, t_\text{surf}, t_\text{scat}\},
\label{eq:event}
\end{equation}
y su tipo se determina por el índice del mínimo, garantizando exclusividad mutua.
Para la pared circular, $t_\text{wall}$ es la raíz positiva de la ecuación
cuadrática de intersección entre el rayo y el cilindro de radio $r_\text{env}$.

\src{Intersección con pared circular: \code{simulation\_engine.py}, líneas
488--501; con caja rectangular: líneas 502--509; muestreo de camino libre y
clasificación de eventos: líneas 525--535.}

\subsection{Contabilización en los planos (\emph{tally})}

Antes de procesar el evento, se detecta si el segmento de rayo
$[\,z_0,\ z_0+t_\text{event}\,d_z\,]$ cruza alguna profundidad objetivo $z^{*}$,
mediante el cambio de signo $(z_0-z^{*})(z_\text{end}-z^{*})<0$. En el punto de cruce
se registra el peso $W$ (con la ponderación pineal de \eqref{eq:pineal} si
corresponde). Esta contabilización por cruce, en ambas direcciones, equivale al
flujo radiante total que atraviesa el plano por unidad de área, $|E_d|+|E_u|$, y se
reduce a la irradiancia direccional clásica cuando no hay dispersión.

\src{\code{simulation\_engine.py}, líneas 537--565.}

\subsection{Función de fase de Henyey--Greenstein}

Ante un evento de dispersión, la nueva dirección se muestrea de la función de fase
de Henyey--Greenstein, parametrizada por el factor de asimetría $g\in(-1,1)$. El
coseno del ángulo de dispersión se obtiene por inversión analítica de la CDF:

\begin{equation}
\cos\theta =
\begin{dcases}
\frac{1}{2g}\left[\,1+g^2-\left(\frac{1-g^2}{1-g+2g\,\xi_1}\right)^{\!2}\right], & g\neq 0,\\[8pt]
1-2\,\xi_1, & g=0,
\end{dcases}
\label{eq:hg}
\end{equation}

con $\xi_1\sim U(0,1)$ y un ángulo azimutal $\phi=2\pi\xi_2$ uniforme. La dirección
saliente se reconstruye en una base ortonormal anclada a la dirección incidente.
Tras la dispersión, el peso del paquete se multiplica por el albedo de dispersión
simple $\omega$, lo que contabiliza la fracción absorbida en el evento
($1-\omega$).

\src{\fn{sample\_henyey\_greenstein}, \code{simulation\_engine.py}; aplicación y
atenuación por $\omega$ en el evento de dispersión interna del bucle Monte Carlo.}

\subsection{Función de fase de Fournier--Forand (alta fidelidad)}\label{sec:ff}

La función de fase de Henyey--Greenstein, aunque conveniente, representa pobremente
el lóbulo de retrodispersión y ata la fracción retrodispersada a un único parámetro
$g$. Para aplicaciones más exigentes, el simulador implementa como alternativa
seleccionable la función de fase de Fournier--Forand
(\code{phase\_function='fournier\_forand'})~\cite{fournier,mobley}, que reproduce a
la vez el pico difractivo hacia adelante y el lóbulo de retrodispersión a partir de
dos parámetros físicos: el índice de refracción real de las partículas $n$ (relativo
al agua) y la pendiente de Junge $\mu$ de la distribución de tamaños:

\begin{equation}
\tilde\beta(\theta) =
\frac{\nu(1-\delta)-(1-\delta^{\nu}) + \bigl[\delta(1-\delta^{\nu})-\nu(1-\delta)\bigr]\sin^{-2}(\theta/2)}
{4\pi\,(1-\delta)^2\,\delta^{\nu}}
+ \frac{1-\delta_{180}^{\nu}}{16\pi\,(\delta_{180}-1)\,\delta_{180}^{\nu}}\,(3\cos^2\theta - 1),
\label{eq:ff}
\end{equation}

con $\nu=(3-\mu)/2$, $\delta = \tfrac{4}{3(n-1)^2}\sin^2(\theta/2)$ y
$\delta_{180}=\delta(\pi)$. Su fracción de retrodispersión tiene forma analítica
(integral de $90^\circ$ a $180^\circ$):

\begin{equation}
\frac{b_b}{b} = 1 - \frac{1-\delta_{90}^{\,\nu+1}-\tfrac12\,(1-\delta_{90}^{\nu})}{(1-\delta_{90})\,\delta_{90}^{\nu}},
\qquad \delta_{90}=\frac{2}{3(n-1)^2}.
\label{eq:ffbb}
\end{equation}

\textbf{Desacople de la retrodispersión.} En lugar de exigir al usuario un índice de
refracción, el simulador toma como entrada la razón de retrodispersión objetivo
$b_b/b$ y resuelve por bisección el $n$ que la satisface a $\mu$ fija
(\eqref{eq:ffbb} es monótona en $n$). Así $b_b$ queda \emph{desacoplada} de la
asimetría, a diferencia de Henyey--Greenstein. Como la fase de Fournier--Forand no
admite inversión analítica de su CDF, las direcciones se muestrean por inversión
numérica de la CDF angular tabulada (densa en el pico hacia adelante), construida una
sola vez por simulación ya que depende sólo de $(n,\mu)$.

\src{\fn{fournier\_forand\_phase}, \fn{ff\_backscatter\_fraction},
\fn{ff\_n\_from\_backscatter}, \fn{build\_ff\_inverse\_cdf} y
\fn{sample\_fournier\_forand} en \code{simulation\_engine.py}; selección de fase y
construcción de la CDF inversa al inicio de \fn{SimulationEngine.run}.}

\nota{Validación numérica: la fracción de retrodispersión muestreada reproduce el
objetivo $b_b/b$ con error $\lesssim 0{,}0005$ en el rango $0{,}005$--$0{,}10$
(oceánico a costero), y el albedo de asimetría efectivo $\langle\cos\theta\rangle$
decrece coherentemente al aumentar $b_b/b$. Por defecto, si no se especifica
$b_b/b$, se usa el valor equivalente de Henyey--Greenstein para la misma $g$, de modo
que el cambio de fase es continuo.}

\subsection{Condiciones de frontera}

\begin{description}[style=nextline,leftmargin=1.4em]
  \item[Pared lateral.] Reflexión \emph{lambertiana}: la nueva dirección se muestrea
  con densidad proporcional al coseno respecto a la normal interior, y el peso se
  multiplica por la reflectancia de pared $r_\text{wall}$ (por defecto $0{,}15$). El
  muestreo coseno-ponderado usa $\cos\theta=\sqrt{\xi_1}$.
  \item[Suelo.] Absorción total: el peso se anula ($W\leftarrow0$).
  \item[Superficie.] Fresnel con posible reflexión interna total, como en
  \S\ref{sec:fresnel}: el peso se multiplica por la reflectancia $R$ y la dirección
  se refleja especularmente.
\end{description}

\src{\fn{sample\_lambertian} (líneas 78--94) y manejo de fronteras: pared (líneas
567--588), suelo (590--592), superficie (594--615) en
\code{simulation\_engine.py}.}

\subsection{Terminación insesgada: ruleta rusa}

Para acotar el costo computacional sin introducir sesgo, los paquetes de bajo peso
se terminan estocásticamente mediante ruleta rusa a partir del rebote
$n_\text{rr}=4$ (máximo $20$ rebotes). Un paquete con peso $W$ por debajo de un
umbral $W_\text{th}$ sobrevive con probabilidad
$p = \min(W/W_\text{th},\,1)$ acotada inferiormente; los supervivientes se
reescalan $W\leftarrow W/p$, preservando la esperanza del estimador:

\begin{equation}
\mathbb{E}[W_\text{post}] = p\cdot\frac{W}{p} + (1-p)\cdot 0 = W.
\label{eq:rr}
\end{equation}

El umbral $W_\text{th}$ se fija en una fracción del peso inicial mediano de los
paquetes.

\src{Ruleta rusa (``CORRECCIÓN \#6''): \code{simulation\_engine.py}, líneas
467--473 y 626--638.}

% =====================================================
\section{Agregación espacial y salida}\label{sec:agregacion}
% =====================================================

Los cruces registrados en cada plano objetivo se agregan en una malla regular de
$100\times100$ celdas que cubre la sección del recinto. La irradiancia en cada celda
es la suma de los pesos de los rayos que la cruzan dividida por el área de la celda:

\begin{equation}
E_{ij} \;=\; \frac{1}{\Delta A}\sum_{k\,\in\,\text{celda }(i,j)} \Phi_k,
\qquad \Delta A = \Delta x\,\Delta y.
\label{eq:bin}
\end{equation}

Sobre esta malla se calculan los estadísticos por capa (promedio, mínimo, máximo),
el área iluminada por encima de un umbral de contorno, el flujo total
$\sum E_{ij}\Delta A$ y, opcionalmente, el recorte a una región de interés (ROI)
cilíndrica o paralelepipédica. Los factores medios $f_\text{lux}$ y $f_\text{ppfd}$
(\S\ref{sec:unidades}) convierten la irradiancia radiométrica a iluminancia y PPFD.
El perfil de profundidad agrega los promedios por capa ponderados por volumen.

\src{Malla y agregación: \code{app\_sim.py}, líneas 420--467; estadísticos de capa
y ROI: líneas 469--530.}

% =====================================================
\section{Calidad de la luz: ángulo de matiz}\label{sec:calidad}
% =====================================================

Además de la \emph{intensidad} (irradiancia, lux, PPFD), el simulador reporta la
\emph{calidad} espectral de la luz —su color— mediante el ángulo de matiz de la
irradiancia descendente $\alpha_E$, siguiendo a Lee et al.\ (2022)~\cite{lee2022}.
Esto aprovecha el carácter espectral y tridimensional del trazador: cada rayo porta
su longitud de onda, de modo que en cualquier plano o celda puede reconstruirse el
espectro local $E_d(\lambda)$ y resumir su color en un único escalar objetivo.

A partir de las contribuciones por rayo se computan los triestímulos CIE con las
funciones de igualación de color del observador estándar de 1931
$\bar x,\bar y,\bar z$:

\begin{equation}
X = \sum_k \Phi_k\,\bar x(\lambda_k),\quad
Y = \sum_k \Phi_k\,\bar y(\lambda_k),\quad
Z = \sum_k \Phi_k\,\bar z(\lambda_k),
\label{eq:tristimulus}
\end{equation}

las coordenadas de cromaticidad $x=X/(X{+}Y{+}Z)$, $y=Y/(X{+}Y{+}Z)$, y el ángulo de
matiz medido respecto al punto blanco equienergético $E=(1/3,1/3)$:

\begin{equation}
\alpha_E = \operatorname{atan2}\!\bigl(y-\tfrac13,\; x-\tfrac13\bigr) \ \in [0,360^\circ),
\qquad
s = \sqrt{(x-\tfrac13)^2 + (y-\tfrac13)^2},
\label{eq:hueangle}
\end{equation}

donde $s$ es la saturación (distancia radial al blanco), que cuantifica la fiabilidad
del matiz: para $s\to 0$ el color es casi neutro y $\alpha_E$ pierde significado. El
simulador entrega $\alpha_E$ por capa y por región de interés, y un mapa de matiz por
celda (tres histogramas ponderados por las funciones de igualación), revelando la
variación 3-D del color de la luz.

\src{\fn{cie\_cmf}, \fn{hue\_angle\_from\_xyz}, \fn{hue\_angle\_from\_spectrum} en
\code{simulation\_engine.py}; tristímulos por celda, $\alpha_E$ por capa/ROI y mapa
de matiz en \code{app\_sim.py}; figura en \fn{plotter.plot\_hue\_angle\_heatmap}.}

\nota{Convención: rojo $\approx 0^\circ/360^\circ$, verde $\approx 110^\circ$, azul
$\approx 240^\circ$. La dirección del cambio con la profundidad depende del agua: en
agua clara el rojo se absorbe primero y la luz \emph{azulea} ($\alpha_E$ crece); en
agua rica en CDOM el azul se absorbe y la luz vira a verde/rojo ($\alpha_E$
decrece). La capa de cambio rápido es la \emph{cromoclina} de Lee et al.\ (2022). El
valor debe interpretarse junto con la saturación $s$: matices con $s$ muy bajo o pocos
impactos por celda son poco fiables. La igualación de color usa la tabla CIE 1931 a
\SI{10}{\nano\meter}, suficiente para un indicador de matiz.}

% =====================================================
\section{Profundidad de Secchi equivalente}\label{sec:secchi}
% =====================================================

Como indicador sintético de la transparencia del agua, el simulador reporta una
profundidad de disco de Secchi equivalente $Z_\text{SD}$. Esta magnitud
\emph{no} se obtiene del trazado de rayos, sino que se deriva analíticamente de los
coeficientes ópticos del escenario, de modo que el usuario disponga de una
referencia intuitiva y comparable con observaciones de campo.

\subsection{Modelo acoplado de Preisendorfer}

La teoría de la visibilidad del disco de Secchi establece que su profundidad de
desaparición es inversamente proporcional a la \emph{suma} del coeficiente de
atenuación de haz $c$ (que gobierna el contraste de la imagen del disco a lo largo
de la línea de visión) y el coeficiente de atenuación difusa $K_d$ (que gobierna la
iluminación del campo de fondo). Siguiendo la formulación clásica acoplada de
Preisendorfer (1986)~\cite{preisendorfer}, el simulador adopta

\begin{equation}
Z_\text{SD} \;\approx\; \frac{\Gamma}{\,c + K_d\,}, \qquad \Gamma \approx 8{,}69,
\label{eq:secchi}
\end{equation}

donde la constante de acoplamiento $\Gamma\approx8{,}69$ agrupa el umbral de
contraste visual del observador y la geometría de iluminación. Esta forma es
preferible a las relaciones de un solo coeficiente cuando se dispone tanto de $c$
como de $K_d$, porque captura explícitamente que la pérdida de contraste del disco
está dominada por la atenuación de haz, no por la difusa.

\src{\fn{\_secchi\_preisendorfer} y lógica de selección:
\code{app\_sim.py}, líneas 644--663. La constante y la cita a Preisendorfer (1986)
están en los comentarios del código (líneas 644--647).}

\subsection{Estimación cruzada cuando falta un coeficiente}

En la práctica, según el modo óptico activo, el simulador puede conocer sólo uno de
los dos coeficientes. Para garantizar que una misma agua entregue el mismo Secchi
con independencia de si se ingresa $c$ o $K_d$, el cierre es \emph{simétrico} y la
fórmula acoplada \eqref{eq:secchi} se aplica de forma unificada a ambos casos:

\begin{enumerate}[itemsep=3pt]
  \item \textbf{Sólo $c$ conocido} (\code{atten\_coef\_type='c'} en \code{kd\_fijo} o
  \code{kd\_espectral}, o \code{scattering} escalar): se estima $K_d$ con el cierre
  bio-óptico de \eqref{eq:cfromkd}, fijando $\omega=0{,}8$, $g=0{,}85$ y
  $\bar\mu_d=0{,}85$,
  \begin{equation}
  K_d \;\approx\; c\,\frac{1-\omega g}{\bar\mu_d}.
  \label{eq:kdfromc}
  \end{equation}
  \item \textbf{Sólo $K_d$ conocido} (\code{atten\_coef\_type='kd'}): se estima $c$
  con la relación inversa de \eqref{eq:cfromkd},
  \begin{equation}
  c \;\approx\; K_d\,\frac{\bar\mu_d}{1-\omega g}.
  \label{eq:cfromkd2}
  \end{equation}
\end{enumerate}

En ambos casos se evalúa luego \eqref{eq:secchi} con el par $(c, K_d)$ completado.
Como $K_d$ y $c$ quedan ligados por el mismo factor en ambas direcciones, el
resultado es idéntico para una misma agua física: por ejemplo $K_d=0{,}2\;\si{\per\meter}$
y su $c$ equivalente $\approx0{,}53\;\si{\per\meter}$ dan ambos
$Z_\text{SD}\approx 11{,}9\;\si{\meter}$.

\src{Cierre simétrico unificado y evaluación en la ventana transparente:
\code{app\_sim.py}, función \fn{\_spectral\_kd\_c\_secchi} y bloque
``Profundidad de disco de Secchi equivalente''. El factor $(1-\omega g)/\bar\mu_d$ y
su inverso son los de \eqref{eq:cfromkd}.}

\nota{La relación de Poole--Atkins $Z_\text{SD}\approx 1{,}7/K_d$~\cite{poole} está
disponible como tercer modelo de salida seleccionable (\S\ref{sec:secchimodelos}),
además de usarse para la \emph{ingesta} de datos satelitales (conversión inversa
$K_d\leftarrow Z_\text{SD}$ en \code{optical\_lookup.py}). La diferencia clave con la
implementación anterior es que ahora se aplica de forma coherente: deriva $K_d$ desde
$c$ con el mismo cierre cuando hace falta, de modo que una misma agua no arroja dos
profundidades distintas según el coeficiente ingresado.}

\subsection{Teoría revisada de Lee et al.\ (2015)}\label{sec:lee2015}

Lee et al.~\cite{lee2015} revisan críticamente la teoría clásica de visibilidad
sobre la que se construye \eqref{eq:secchi}. Su argumento central es que esa teoría
trata al disco de Secchi como un objeto puntual e ignora la alta resolución angular
del ojo humano ($\sim$\SI{0.5}{\arcmin}); como un disco de $\sim$\SI{30}{\centi\meter}
sigue siendo un objeto extenso a decenas de metros, la Ley de Reducción de Contraste
no aplica, y la atenuación del contraste deja de seguir el coeficiente de haz $c$.
Empíricamente, no existe relación universal entre $Z_\text{SD}$ y $c$, mientras que
$Z_\text{SD}$ correlaciona igual o mejor con $K_d$. El modelo mecanístico que
proponen y validan ($N=338$, $R^2=0{,}96$, $\sim$18\,\% de error medio) hace que la
transparencia quede gobernada por el coeficiente de atenuación difusa \emph{mínimo}
del visible —la ventana transparente— y no por $c+K_d$:

\begin{equation}
Z_\text{SD} \;=\; \frac{1}{2{,}5\,K_d^{tr}}\,
\ln\!\left(\frac{|r_T - r_w|}{C_t^{\,r}}\right),
\qquad K_d^{tr} = \min_{\lambda\in[400,700]} K_d(\lambda),
\label{eq:secchilee}
\end{equation}

donde $r_T$ es la reflectancia de radiancia del disco blanco
($R_T=0{,}85$ lambertiano, $r_T=R_T/\pi$), $r_w$ la reflectancia de fondo del agua
en la ventana transparente, y $C_t^{\,r}$ el contraste umbral del ojo
(por defecto $0{,}013\;\si{\per\steradian}$).

El simulador implementa tres modelos como opción seleccionable\label{sec:secchimodelos}
(\code{secchi\_model} $\in\{$\code{preisendorfer}, \code{poole\_atkins}, \code{lee2015}$\}$),
todos con cierre coherente $c\leftrightarrow K_d$ (\S\ref{sec:secchi}). El $K_d(\lambda)$
espectral se obtiene del modo óptico activo (del bio-óptico vía \eqref{eq:kdfromiop}
en modo \code{bio}; de la tabla en \code{kd\_espectral}; o del cierre \eqref{eq:cfromkd}
cuando sólo se conoce $c$), y se toma su mínimo como $K_d^{tr}$. La retrodispersión
entra al modelo de Lee \eqref{eq:secchilee} por \emph{dos} vías: a través de $K_d$
—dominante, cuando se usa el cierre de Lee 2005 \eqref{eq:kdlee} con $b_b$ explícito—
y a través de $r_w$, la reflectancia de fondo dentro del logaritmo de contraste. En
\emph{todos} los modos ópticos, $r_w$ se estima desde las IOP en la ventana
transparente con la relación de Gordon $R(0^-)\approx f\,b_b/(a+b_b)$,
$f\approx0{,}33$, $r_w=R(0^-)/\pi$, usando la retrodispersión activa $b_b=B\,b$ con
$B$ de la fase seleccionada: la fracción analítica de Henyey--Greenstein,

\begin{equation}
B(g) \;=\; \frac{1-g}{2g}\left(\frac{1+g}{\sqrt{1+g^2}}-1\right),
\qquad B(0{,}85)\approx 0{,}036,
\label{eq:hgbackscatter}
\end{equation}

o directamente la razón objetivo $b_b/b$ cuando se usa la fase de Fournier--Forand.
El efecto de $b_b$ sobre $r_w$ es secundario frente al que ejerce vía $K_d$ (porque
$r_T\approx0{,}27$ domina el contraste), pero crece en aguas muy dispersoras y es
físicamente correcto incluirlo.

\src{Primitivas de física: \fn{secchi\_lee2015}, \fn{secchi\_preisendorfer},
\fn{secchi\_poole\_atkins}, \fn{hg\_backscatter\_fraction},
\fn{subsurface\_reflectance} en \code{simulation\_engine.py} (sección ``Profundidad de disco de Secchi
equivalente''). Orquestación, selección de modelo y evaluación espectral por modo:
\code{app\_sim.py}, bloque ``Profundidad de disco de Secchi equivalente''.}

La Tabla~\ref{tab:secchicomp} contrasta ambos modelos evaluados con el modelo
bio-óptico del simulador en tres escenarios de fiordo. La razón $c/K_d^{tr}$ se sitúa
en $2{,}8$--$3{,}8$, dentro del rango ``$c$ es 2--5$\times K_d$'' reportado por
Lee et al.; el cierre inverso del simulador (\S\ref{sec:biooptico}) es por tanto
consistente con esa observación empírica.

\begin{table}[h]
\centering
\caption{Comparación de modelos de Secchi sobre el modelo bio-óptico del simulador
(ventana transparente $\lambda_{tr}$). Valores ilustrativos.}
\label{tab:secchicomp}
\small
\begin{tabular}{lcccccc}
\toprule
Escenario & $\lambda_{tr}$ [\si{\nano\meter}] & $K_d^{tr}$ [\si{\per\meter}] & $c$ [\si{\per\meter}] & $c/K_d^{tr}$ & $Z_\text{SD}^\text{Preis}$ [\si{\meter}] & $Z_\text{SD}^\text{Lee}$ [\si{\meter}] \\
\midrule
Claro  & 550 & 0,27 & 0,76 & 2,8 & 8,5 & 4,5 \\
Típico & 550 & 0,75 & 2,75 & 3,7 & 2,5 & 1,6 \\
Turbio & 635 & 1,50 & 5,68 & 3,8 & 1,2 & 0,8 \\
\bottomrule
\end{tabular}
\end{table}

\nota{Las constantes $\Gamma=8{,}69$ (Preisendorfer), $1{,}7$ (Poole--Atkins) y el
contraste umbral $C_t^{\,r}$ (Lee) son valores nominales de literatura no calibrados
al observador ni al sitio. El modelo de Lee et al.\ (2015) es físicamente más
defendible para discos de Secchi y predice una transparencia gobernada por $K_d$, no
por $c$; en aguas de fiordo/jaula, donde $c\gg K_d$, ambos modelos difieren de forma
material (Tabla~\ref{tab:secchicomp}). El valor reportado debe interpretarse como
indicador comparativo de transparencia, no como sustituto de una lectura de disco en
terreno. Para mayor fidelidad espectral de $K_d$, Lee et al.\ emplean QAA~\cite{qaa}
y el modelo de $K_d$ de Lee et al.~\cite{lee2005}, más elaborados que el cierre tipo
Kirk del simulador.}

% =====================================================
\section{Presets bio-ópticos desde teledetección}\label{sec:presets}
% =====================================================

El módulo \code{optical\_lookup.py} genera presets de turbidez (\emph{claro},
\emph{típico}, \emph{turbio}) compatibles con el modo \code{bio}, a partir de
observaciones satelitales o proxy por centro de cultivo. El objetivo es transformar
una serie temporal heterogénea de productos de teledetección (con distintas fuentes,
variables y niveles de cobertura) en tres tríos $(\text{TSS}, a_\text{CDOM}(440),
\text{Chl})$ que alimenten directamente el modelo bio-óptico de
\S\ref{sec:biooptico}. El proceso consta de cuatro etapas: (i) ingesta y
homogeneización de observaciones, (ii) conversión de variables proxy, (iii)
agregación estadística por escenario, y (iv) cuantificación de la confianza.

\subsection{Ingesta y conversiones proxy}\label{sec:proxy}

Las observaciones se leen de un CSV con columnas para TSS, SPM, turbidez en FNU,
Chl-a, CDOM (a 440 o 443~\si{\nano\meter}), $K_{d,490}$, profundidad de Secchi
$Z_\text{SD}$ y metadatos de calidad. Cuando una variable primaria falta, se aplica
una jerarquía de sustituciones proxy, cada una con su trazabilidad registrada:

\begin{description}[style=nextline,leftmargin=1.4em]
  \item[TSS desde SPM.] Si no hay TSS pero sí materia particulada en suspensión
  (SPM), se usa directamente $\text{TSS}=\text{SPM}$ (equivalencia operacional).
  \item[TSS desde turbidez FNU.] Si tampoco hay SPM, se convierte la turbidez
  nefelométrica mediante una calibración lineal local
  \begin{equation}
  \text{TSS}\,[\si{\milli\gram\per\liter}] \;=\; m\cdot \text{FNU} + b,
  \label{eq:fnu}
  \end{equation}
  con pendiente $m$ e intercepto $b$ configurables. Por defecto $m=1$, $b=0$, es
  decir la equivalencia conservadora $1\,\text{FNU}\approx\SI{1}{\milli\gram\per\liter}$,
  que el código declara explícitamente como ``calibración local requerida''. La
  relación FNU$\to$TSS es lineal en el rango de validez de la forma semianalítica de
  Nechad et al.~\cite{nechad} para sólidos en suspensión.
  \item[CDOM 443 $\to$ 440.] Si sólo se dispone de $a_\text{CDOM}(443)$, se traslada
  a la referencia de \SI{440}{\nano\meter} con la pendiente espectral típica
  $S_\text{CDOM}=\SI{0.015}{\per\nano\meter}$:
  \begin{equation}
  a_\text{CDOM}(440) \;=\; a_\text{CDOM}(443)\,e^{\,S_\text{CDOM}\,(443-440)}
  \;=\; a_\text{CDOM}(443)\,e^{\,0{,}015\times 3}.
  \label{eq:cdomshift}
  \end{equation}
  \item[$K_{d,490}$ desde Secchi.] Si falta $K_{d,490}$ pero existe $Z_\text{SD}$,
  se estima por la relación de Poole--Atkins~\cite{poole}
  $K_{d,490}\approx 1{,}7/Z_\text{SD}$ (relación de Poole--Atkins~\cite{poole}).
\end{description}

\src{Ingesta y proxies: \code{optical\_lookup.py}, líneas 26--99
(\fn{load\_observations}, \fn{\_normalize\_proxy\_observations}). Calibración
FNU$\to$TSS: líneas 96--108 (\fn{\_tss\_from\_turbidity\_fnu}). Conversión CDOM:
líneas 55--58. Estimación $K_d$ desde Secchi: líneas 221--222.}

\subsection{Estimador bio-óptico de $K_{d,490}$}\label{sec:kd490}

Para coherencia interna, cuando se necesita un $K_{d,490}$ a partir de las
concentraciones, el módulo evalúa el mismo modelo de cuatro componentes de
\S\ref{sec:biooptico} fijado en \SI{490}{\nano\meter} y lo cierra con la relación de
Kirk \eqref{eq:kdfromiop}:

\begin{equation}
K_{d,490} \;=\; \frac{\bigl[a_w + a_\text{CDOM}(440)e^{-S_\text{CDOM}(490-440)} + a_\phi^{*}\,\text{Chl}\bigr] + (1-g)\,b^{*}\,\text{TSS}}{\bar\mu_d},
\label{eq:kd490}
\end{equation}

con los coeficientes evaluados en \SI{490}{\nano\meter}: $a_w=0{,}026$,
$a_\phi^{*}=0{,}012$, $b^{*}=0{,}35$, y $g=\bar\mu_d=0{,}85$. Este estimador permite
``ajustar'' las concentraciones por defecto de una clase de agua para que reproduzcan
un $K_{d,490}$ observado: se calcula la razón
$r=K_{d,490}^\text{obs}/K_{d,490}^\text{default}$, se acota a $r\in[0{,}35,\,3{,}0]$
para evitar extrapolaciones, y se escalan TSS y CDOM por $r$.

\src{\fn{\_estimate\_kd490}: \code{optical\_lookup.py}, líneas 336--346.
\fn{\_fit\_defaults\_to\_kd} con el factor de escala acotado: líneas 349--359.}

\subsection{Agregación por escenario y cuantiles}\label{sec:cuantiles}

Los tres escenarios se definen como cuantiles de la distribución empírica de cada
variable sobre las observaciones que coinciden con el centro:

\begin{equation}
\text{claro} \to q_{25},\qquad \text{típico} \to q_{50}\ (\text{mediana}),\qquad \text{turbio} \to q_{75}.
\label{eq:cuantiles}
\end{equation}

Para cada escenario y variable, si existe el cuantil observado se usa éste; si no,
se recurre al valor por defecto de la clase de agua ajustado por $K_d$
(\S\ref{sec:kd490}). El cuantil empírico se calcula por interpolación lineal entre
órdenes, robusto a tamaños de muestra pequeños.

\src{\fn{\_scenario\_values}: \code{optical\_lookup.py}, líneas 362--391; función de
cuantil interpolado \fn{\_quantile}: líneas 77--88.}

\subsection{Climatología estacional por semana ISO}\label{sec:isoweek}

El perfil estacional usa la semana ISO en lugar de fechas arbitrarias, lo que alinea
años distintos sobre un calendario común de 53 semanas. La agregación es jerárquica
y con ponderación igual por año, para que un año con mayor cobertura satelital no
domine el clima:

\begin{equation}
\bar{x}_w \;=\; \frac{1}{|Y_w|}\sum_{y\in Y_w} \operatorname{med}\bigl(\{x_{w,y,d}\}_d\bigr),
\label{eq:isoweek}
\end{equation}

donde primero se toma la mediana de todos los días $d$ dentro de la semana $w$ del
año $y$, y luego se promedian los $|Y_w|$ años disponibles con igual peso. Una
semana se considera útil cuando reúne al menos cuatro días válidos repartidos en dos
o más años. El endpoint \code{/api/optical\_weekly\_profile} devuelve las 53 semanas
con su cobertura, medianas, rangos intercuartílicos y los presets
\emph{claro}/\emph{típico}/\emph{turbio}.

\src{\fn{\_aggregate\_week\_rows\_by\_year} (mediana intra-año) y combinación
inter-año con peso igual: \code{optical\_lookup.py}, líneas 131--160. Asignación de
semana ISO: \fn{\_observation\_iso\_week}, líneas 118--128.}

\subsection{Cuantificación de la confianza}\label{sec:confianza}

Cada juego de presets se acompaña de un nivel de confianza derivado del volumen y
diversidad temporal de las observaciones: \emph{baja} (sin datos, sólo proxies por
clase de agua), \emph{baja-media} ($<4$ observaciones), \emph{media}
($\geq 4$), y \emph{media-alta} ($\geq 10$ observaciones en $\geq 5$ días
distintos). Se reportan además medianas y rangos intercuartílicos (IQR) de TSS, Chl,
CDOM, $K_{d,490}$ y turbidez FNU, la fracción de TSS derivada por proxy, las fuentes
de observación y, cuando existen, las incertidumbres porcentuales por píxel y el
número de píxeles válidos.

\src{Construcción del diccionario de confianza: \code{optical\_lookup.py}, líneas
384--457.}

\subsection{Conectores satelitales}\label{sec:conectores}

Los conectores remotos se desacoplan en el paquete \code{optical\_sources/} y se
seleccionan en modo \code{auto} según el tipo de centro. La
Tabla~\ref{tab:conectores} resume sus productos.

\begin{table}[h]
\centering
\caption{Conectores de teledetección implementados (\code{optical\_sources/}).}
\label{tab:conectores}
\small
\begin{tabular}{p{2.6cm}p{3.2cm}p{3.0cm}p{3.4cm}}
\toprule
\textbf{Conector} & \textbf{Producto / dataset} & \textbf{Variables} & \textbf{Acceso} \\
\midrule
NOAA CoastWatch & DINEOF global diario (S-NPP/N20 VIIRS) & \code{chlor\_a}, \code{kd\_490} & ERDDAP público, sin credenciales \\
Copernicus Marine & GlobColour L3 multi 4\,km diario (\code{my}/\code{nrt}) & CHL, KD490, SPM, CDM + incertidumbres & API CMEMS \\
NASA OceanColor & VIIRS-N y PACE-OCI L3m 4\,km diario & \code{chlor\_a}, \code{Kd\_490}, \code{adg\_443} & Descarga L3m con caché local \\
Sentinel-2 & Salidas ACOLITE (\code{.nc}/\code{.csv}) & turbidez FNU, \code{rhow\_665} & Directorio local ACOLITE/DSF \\
\bottomrule
\end{tabular}
\end{table}

\src{\code{noaa\_coastwatch.py}, \code{copernicus.py}, \code{nasa\_oceancolor.py},
\code{sentinel2.py}. Mapeo de variables NASA: \code{nasa\_oceancolor.py}, líneas
14--19 y 158. Datasets Copernicus: \code{copernicus.py}, líneas 15--23 y 162--163.}

\paragraph{Cadena Sentinel-2 / ACOLITE.}
La ruta Sentinel-2 mantiene separadas tres capas físicas independientes, lo que
permite auditar y recalibrar cada una por separado: (1)~la corrección atmosférica
ACOLITE/DSF que produce la reflectancia de agua $\rho_w$; (2)~la estimación de
turbidez; y (3)~la calibración local FNU$\to$TSS de \eqref{eq:fnu}. Cuando ACOLITE
ya entrega un producto de turbidez, el conector lo usa directamente; cuando sólo hay
reflectancia $\rho_w(665)$, aplica la forma semianalítica de
Nechad et al.~\cite{nechad,vanhellemont}:

\begin{equation}
T \;=\; \frac{A_T\,\rho_w}{\,1 - \rho_w/C\,} + B_T,
\label{eq:nechad}
\end{equation}

con los coeficientes $A_T$, $C$ y (opcionalmente) $B_T$ configurables por variables
de entorno (\code{SENTINEL2\_NECHAD\_AT}, \code{\_C}, \code{\_BT}). La forma es válida
sólo para $0<\rho_w<C$. A las observaciones Sentinel-2 sin incertidumbre explícita se
les asigna por defecto la RMSE de validación local
$\SI{0.66}{\text{FNU}}$ (etiqueta \code{RELONCAVI\_NV09}).

\src{Forma de Nechad \fn{\_nechad\_turbidity}: \code{sentinel2.py}, líneas 62--67;
coeficientes por entorno: líneas 47--60. RMSE por defecto
\code{RELONCAVI\_NV09\_RMSE\_FNU}: \code{optical\_lookup.py}, líneas 22 y 97--98.}

\paragraph{Selección geográfica.}
Cada consulta define una caja geográfica alrededor de las coordenadas del centro a
partir de un radio de \emph{buffer} en metros, convertido a grados con la corrección
del coseno de la latitud para la longitud:
$\Delta\text{lon} = \Delta\text{lat}/\max(\cos\phi_\text{lat},\,0{,}2)$. Los píxeles
dentro de la caja se agregan por mediana.

\src{\fn{\_buffer\_to\_box}: \code{sentinel2.py}, líneas 36--39.}

\nota{Las clases de agua por defecto (\code{fjord\_clear}, \code{fjord\_typical},
\code{fjord\_turbid}, \code{coastal\_turbid}) son puntos de partida conservadores
para centros marinos chilenos sin extracción satelital cargada; explícitamente no
son mediciones de campo. La calibración FNU$\to$TSS por defecto ($1{,}0$, $0{,}0$) y
los coeficientes de Nechad deben ajustarse con datos locales antes de usar los
presets con fines cuantitativos. La cadena Sentinel-2 mantiene separadas las tres
capas (corrección atmosférica, turbidez, FNU$\to$TSS) precisamente para no mezclar
sus incertidumbres.}

% =====================================================
\section{Supuestos y límites de validez}\label{sec:limites}
% =====================================================

\begin{itemize}[itemsep=3pt]
  \item \textbf{Estado estacionario y medio no emisor.} No se modela emisión
  termal ni inelástica (fluorescencia, Raman). La luz proviene únicamente de las
  luminarias.
  \item \textbf{Polarización ignorada.} Fresnel se evalúa para luz no polarizada
  \eqref{eq:fresnel}; no se transporta estado de polarización.
  \item \textbf{Superficie plana.} La interfaz aire--agua se trata como un plano
  horizontal: no se modela oleaje ni enfoque por ondas capilares.
  \item \textbf{Fase de dispersión.} La dispersión volumétrica admite dos fases
  seleccionables: Henyey--Greenstein (una $g$) y Fournier--Forand con $b_b/b$
  desacoplada (\S\ref{sec:ff}). Ambas usan una fase única, sin dependencia espectral
  de la asimetría salvo la implícita en $\omega(\lambda)$ del modo \code{bio}.
  \item \textbf{Cierre IOP$\to$AOP aproximado.} El cierre de $K_d$ admite las
  opciones Kirk \eqref{eq:kdfromiop} (con $g$, $\bar\mu_d$ fijos) y Lee et al.\ 2005
  \eqref{eq:kdlee} (con $b_b$ explícito); ambas son parametrizaciones, no soluciones
  exactas de la RTE.
  \item \textbf{Modo RAS con magnitud no calibrada.} El submodelo
  \code{ras\_bardsnes} (\S\ref{sec:ras}) está operativo y reproduce las
  \emph{formas} espectrales de RAS ($\eta_p$, $S_\text{CDOM}$), pero su
  \emph{magnitud} depende de $b^{*}_{550}$ y $\omega_p$, que son calibrables por
  instalación y cuyos valores por defecto no provienen del trabajo original. Sin una
  medida óptica propia del sistema, la ruta es cualitativa en magnitud.
  \item \textbf{Convergencia Monte Carlo.} La incertidumbre estadística decrece como
  $N^{-1/2}$; el número de rayos por defecto es $5\times10^4$. Regiones lejanas o
  muy atenuadas requieren más rayos para reducir el ruido del estimador.
  \item \textbf{Secchi con constantes nominales.} La profundidad de Secchi
  equivalente \eqref{eq:secchi} usa $\Gamma=8{,}69$ y, en su defecto, el factor
  $1{,}7$ de Poole--Atkins; ambos son valores de literatura no calibrados al
  observador ni al sitio, por lo que $Z_\text{SD}$ es un indicador comparativo, no
  una predicción de campo.
  \item \textbf{Calibración satelital pendiente.} Los presets dependen de la
  calibración FNU$\to$TSS (por defecto $1{:}1$) y de los coeficientes de Nechad; sin
  ajuste local son cualitativos. Los productos L3m de NASA usados aquí no incluyen
  una incertidumbre porcentual por píxel equivalente a la de Copernicus.
\end{itemize}

% =====================================================
\section{Nomenclatura}\label{sec:nomenclatura}
% =====================================================

\begin{longtable}{p{2.6cm}p{8.5cm}p{3cm}}
\toprule
\textbf{Símbolo} & \textbf{Descripción} & \textbf{Unidad} \\
\midrule
\endhead
$\Phi$ & Flujo radiante (peso del rayo) & \si{\watt} \\
$I$ & Intensidad radiante & \si{\watt\per\steradian} \\
$E$ & Irradiancia & \si{\watt\per\square\meter} \\
$E_v$ & Iluminancia & \si{\lux} \\
$E_\lambda$ & Irradiancia espectral & \si{\watt\per\square\meter\per\nano\meter} \\
$V(\lambda)$ & Eficiencia luminosa espectral fotópica & --- \\
$K_m$ & Eficacia luminosa máxima ($683$) & \si{\lumen\per\watt} \\
PPFD & Densidad de flujo fotónico fotosintético & \si{\micro\mol\per\square\meter\per\second} \\
$a$ & Coeficiente de absorción & \si{\per\meter} \\
$b$ & Coeficiente de dispersión & \si{\per\meter} \\
$c$ & Coeficiente de atenuación de haz, $c=a+b$ & \si{\per\meter} \\
$K_d$ & Coeficiente de atenuación difusa & \si{\per\meter} \\
$\omega$ & Albedo de dispersión simple, $b/c$ & --- \\
$g$ & Factor de asimetría (Henyey--Greenstein) & --- \\
$\bar\mu_d$ & Coseno medio del campo descendente & --- \\
$S_\text{CDOM}$ & Pendiente espectral del CDOM & \si{\per\nano\meter} \\
$n_1, n_2$ & Índices de refracción (aire, agua) & --- \\
$\theta_i,\theta_t$ & Ángulos de incidencia y transmisión & rad \\
$R, T$ & Reflectancia / transmitancia de Fresnel & --- \\
$r_\text{wall}$ & Reflectancia lambertiana de pared & --- \\
$W$ & Peso del paquete Monte Carlo & \si{\watt} \\
$\text{Chl}$ & Concentración de clorofila-a & \si{\milli\gram\per\cubic\meter} \\
$\text{TSS}$ & Sólidos suspendidos totales & \si{\milli\gram\per\liter} \\
$a_\text{CDOM}(440)$ & Absorción CDOM de referencia a \SI{440}{\nano\meter} & \si{\per\meter} \\
$Z_\text{SD}$ & Profundidad de disco de Secchi (equivalente) & \si{\meter} \\
$\Gamma$ & Constante de acoplamiento de Preisendorfer ($\approx 8{,}69$) & --- \\
$K_d^{tr}$ & $K_d$ mínimo del visible (ventana transparente) & \si{\per\meter} \\
$r_T, r_w$ & Reflectancia de radiancia del disco / del agua de fondo & \si{\per\steradian} \\
$C_t^{\,r}$ & Contraste umbral del ojo (Lee et al.) & \si{\per\steradian} \\
$B(g)$ & Fracción de retrodispersión de Henyey--Greenstein & --- \\
$b_b$ & Coeficiente de retrodispersión, $b_b=B\,b$ & \si{\per\meter} \\
$n$ & Índice de refracción de partícula (Fournier--Forand) & --- \\
$\mu$ & Pendiente de Junge de la distribución de tamaños & --- \\
$\theta_a$ & Ángulo cenital solar en aire (cierre Lee 2005) & \si{\degree} \\
$\alpha_E$ & Ángulo de matiz de la irradiancia descendente & \si{\degree} \\
$\bar x,\bar y,\bar z$ & Funciones de igualación de color CIE 1931 & --- \\
$x,y$ & Coordenadas de cromaticidad CIE & --- \\
$s$ & Saturación (distancia al punto blanco) & --- \\
FNU & Unidad nefelométrica de formazina (turbidez) & FNU \\
$T$ & Turbidez (forma de Nechad) & FNU \\
$\rho_w$ & Reflectancia de agua (corregida por atmósfera) & --- \\
$A_T, B_T, C$ & Coeficientes de la forma de Nechad & --- \\
$q_{25},q_{50},q_{75}$ & Cuantiles de escenario (claro/típico/turbio) & --- \\
\bottomrule
\end{longtable}

% =====================================================
\begin{thebibliography}{99}
\addcontentsline{toc}{section}{Referencias}

\bibitem{mobley} Mobley, C.D. (1994). \emph{Light and Water: Radiative Transfer in
Natural Waters}. Academic Press.

\bibitem{kirk} Kirk, J.T.O. (2011). \emph{Light and Photosynthesis in Aquatic
Ecosystems}, 3rd ed. Cambridge University Press.

\bibitem{hg} Henyey, L.G. \& Greenstein, J.L. (1941). Diffuse radiation in the
galaxy. \emph{The Astrophysical Journal}, 93, 70--83.

\bibitem{fournier} Fournier, G.R. \& Forand, J.L. (1994). Analytic phase function
for ocean water. \emph{Proc. SPIE 2258, Ocean Optics XII}, 194--201.

\bibitem{popefry} Pope, R.M. \& Fry, E.S. (1997). Absorption spectrum (380--700 nm)
of pure water. II. Integrating cavity measurements. \emph{Applied Optics}, 36(33),
8710--8723.

\bibitem{smithbaker} Smith, R.C. \& Baker, K.S. (1981). Optical properties of the
clearest natural waters (200--800 nm). \emph{Applied Optics}, 20(2), 177--184.

\bibitem{bricaud} Bricaud, A., Babin, M., Morel, A. \& Claustre, H. (1995).
Variability in the chlorophyll-specific absorption coefficients of natural
phytoplankton. \emph{Journal of Geophysical Research}, 100(C7), 13321--13332.

\bibitem{bricaud_cdom} Bricaud, A., Morel, A. \& Prieur, L. (1981). Absorption by
dissolved organic matter of the sea (yellow substance) in the UV and visible
domains. \emph{Limnology and Oceanography}, 26(1), 43--53.

\bibitem{babin} Babin, M., Stramski, D., Ferrari, G.M., Claustre, H., Bricaud, A.,
Obolensky, G. \& Hoepffner, N. (2003). Variations in the light absorption
coefficients of phytoplankton, non-algal particles, and dissolved organic matter in
coastal waters around Europe. \emph{Journal of Geophysical Research}, 108(C7), 3211.

\bibitem{morel} Morel, A. \& Maritorena, S. (2001). Bio-optical properties of
oceanic waters: A reappraisal. \emph{Journal of Geophysical Research}, 106(C4),
7163--7180.

\bibitem{ioccg} IOCCG (2006). \emph{Remote Sensing of Inherent Optical Properties:
Fundamentals, Tests of Algorithms, and Applications}. Lee, Z.-P. (Ed.), Reports of
the International Ocean-Colour Coordinating Group, No.~5, IOCCG, Dartmouth.

\bibitem{kirk1984} Kirk, J.T.O. (1984). Dependence of relationship between inherent
and apparent optical properties of water on solar altitude. \emph{Limnology and
Oceanography}, 29(2), 350--356.

\bibitem{lee2005} Lee, Z.-P., Du, K.-P. \& Arnone, R. (2005). A model for the
diffuse attenuation coefficient of downwelling irradiance. \emph{Journal of
Geophysical Research}, 110, C02016.

\bibitem{qaa} Lee, Z.-P., Carder, K.L. \& Arnone, R.A. (2002). Deriving inherent
optical properties from water color: A multiband quasi-analytical algorithm for
optically deep waters. \emph{Applied Optics}, 41(27), 5755--5772.

\bibitem{gordon1975} Gordon, H.R., Brown, O.B. \& Jacobs, M.M. (1975). Computed
relationships between the inherent and apparent optical properties of a flat
homogeneous ocean. \emph{Applied Optics}, 14(2), 417--427.

\bibitem{preisendorfer} Preisendorfer, R.W. (1986). Secchi disk science: Visual
optics of natural waters. \emph{Limnology and Oceanography}, 31(5), 909--926.

\bibitem{poole} Poole, H.H. \& Atkins, W.R.G. (1929). Photo-electric measurements of
submarine illumination throughout the year. \emph{Journal of the Marine Biological
Association of the UK}, 16(1), 297--324.

\bibitem{lee2015} Lee, Z.-P., Shang, S., Hu, C., Du, K., Weidemann, A., Hou, W.,
Lin, J. \& Lin, G. (2015). Secchi disk depth: A new theory and mechanistic model for
underwater visibility. \emph{Remote Sensing of Environment}, 169, 139--149.

\bibitem{lee2018} Lee, Z.-P., Shang, S., Du, K. \& Wei, J. (2018). Resolving the
long-standing puzzles about the observed Secchi depth relationships. \emph{Limnology
and Oceanography}, 63(6), 2321--2336.

\bibitem{lee2022} Lee, Z.-P., Shang, S., Li, Y., Luis, K., Dai, M. \& Wang, Y.
(2022). Three-dimensional variation in light quality in the upper water column
revealed with a single parameter. \emph{IEEE Transactions on Geoscience and Remote
Sensing}, 60, 4203110.

\bibitem{nechad} Nechad, B., Ruddick, K.G. \& Park, Y. (2010). Calibration and
validation of a generic multisensor algorithm for mapping of total suspended matter
in turbid waters. \emph{Remote Sensing of Environment}, 114(4), 854--866.

\bibitem{vanhellemont} Vanhellemont, Q. \& Ruddick, K. (2018). Atmospheric correction
of metre-scale optical satellite data for inland and coastal water applications.
\emph{Remote Sensing of Environment}, 216, 586--597. (ACOLITE / Dark Spectrum
Fitting.)

\bibitem{wyman} Wyman, C., Sloan, P.-P. \& Shirley, P. (2013). Simple analytic
approximations to the CIE XYZ color matching functions. \emph{Journal of Computer
Graphics Techniques}, 2(2), 1--11.

\bibitem{cie} CIE (2018). \emph{Colorimetry}, 4th ed., CIE 015:2018. (Función
$V(\lambda)$ y eficacia luminosa $K_m=683\,$lm/W.)

\bibitem{ies-tm33} Illuminating Engineering Society (2018). \emph{IES TM-33-18:
Standard Format for the Electronic Transfer of Luminaire Optical Data}. IES.

\bibitem{ies-lm63} Illuminating Engineering Society (2019). \emph{IES LM-63-19:
Standard File Format for the Electronic Transfer of Photometric Data}. IES.

\bibitem{bardsnes} Bårdsnes, K. (2020). \emph{Light attenuation in
recirculating aquaculture systems (RAS)}. (De aquí se toman las formas espectrales
$\eta_p\approx1{,}8$ y $S_\text{CDOM}\approx0{,}0141\;\mathrm{nm^{-1}}$ y la
regresión turbidez$\to$TSS de tanque; la magnitud absoluta no es transferible.)

\end{thebibliography}

\end{document}
